\documentclass[10pt,a4paper]{amsart}
\usepackage[margin=19mm]{geometry}
\usepackage{amsmath,amssymb,mathtools}
\usepackage[scr=rsfs,cal=euler]{mathalfa}
\usepackage{xfrac}
\usepackage[unicode,hidelinks,bookmarks=false]{hyperref}
\newcommand{\ie}{i.e.\ }
\newcommand{\cf}{cf.\ }
\newcommand{\resp}{respectively\ }
\newcommand{\Subsection}{\subsection}
\providecommand{\nobreakdash}{\nobreak\hbox{-}}
\setlength{\emergencystretch}{2em}
\multlinegap0pt
\usepackage{bm}
\usepackage{bbm}
\usepackage{dsfont}
\usepackage{xspace}
\usepackage{cancel}
\usepackage{mathtools}
\usepackage{amsthm}
\makeatletter
\@ifundefined{thm}{\newtheorem{thm}{Theorem}}{}
\makeatother
\DeclareMathOperator{\GL}{\mathbf{GL}}
\DeclareMathOperator{\QQ}{\mathbb{Q}}
\DeclareMathOperator{\Sym}{Sym}
\title[Quasiautomorphic forms are isomorphic to vector-valued autom]{Quasiautomorphic forms are isomorphic to vector-valued automorphic forms}
\author{Michael Andrew Henry}
\date{Source: arXiv:2605.21067v1 (CC BY 4.0)}
\begin{document}
\maketitle

\noindent\emph{Source and licence.} Quasiautomorphic forms are isomorphic to vector-valued automorphic forms. Version arXiv:2605.21067v1 (CC BY 4.0), distributed under the Creative Commons Attribution 4.0 International licence, as recorded in the arXiv licence metadata for this exact version and shown on the abstract page https://arxiv.org/abs/2605.21067v1. Full rights notice: licenses/arxiv-2605.21067-CC-BY-4.0.txt.\par\smallskip

\noindent\textit{Editorial scope.} Counted unit: the Thm whose source label is \texttt{None} in the article (source0.tex lines 871--886, source label \texttt{None}), delivered together with the article's own complete original proof of it (source0.tex lines 887--979, one \texttt{proof} environment of 93 lines). No mathematics is written, repaired or re-derived: statement and proof are the article's own text line for line. Required context retained uncounted: none. Every definition, display and label of those lines is retained unchanged. Omitted from this package and not redistributed: 1--870, 980--985. Nothing inside a retained excerpt is omitted or altered: every retained body is delivered exactly as the article's lines are written, so no omission receipt of any kind accompanies this package. Citations: the retained excerpts carry no citation command at all, so no bibliography entry of the article is transcribed and no reference is redistributed. Reference adaptation: none was needed and none was made; every reference inside the delivered unit resolves inside it, so no unresolved reference or citation is produced. Printed slips kept literal: no printed word, symbol, hypothesis or number of the counted statement or of its proof was altered. Layout and class: the article's own class is \texttt{article} with packages including graphicx, svg, tikz-cd, inputenc, amsmath, amsfonts, amssymb, amsthm, booktabs, latexsym, ytableau, tikz, biblatex, float; the wrapper is the lane's pinned \texttt{amsart} editorial wrapper at 10pt on A4, which restates the theorem-like environments the retained text needs and adds the article's own macro definitions that the retained text needs (\texttt{\textbackslash GL}, \texttt{\textbackslash QQ}, \texttt{\textbackslash Sym}), with extra wrapper packages (bm, bbm, dsfont, xspace, cancel, mathtools). The delivered numbering is the unit's own. Assets: no retained span contains a figure environment, a graphics-inclusion command, a PDF-page inclusion, a table or a TikZ picture; no image file is copied into this package, the package declares no asset dependency and no image file is redistributed with it. Licence: version arXiv:2605.21067v1 of this article is distributed under the Creative Commons Attribution 4.0 International licence, as recorded in the arXiv licence metadata for this exact version and shown on the abstract page https://arxiv.org/abs/2605.21067v1. A full rights notice is delivered at licenses/arxiv-2605.21067-CC-BY-4.0.txt. Characters: every retained excerpt is pure ASCII, so the delivered engine, XeLaTeX with its default Latin Modern text font, reports no missing character.
\begin{thm}
    Let $ \mu \geq 3 $ be an integer, let $\zeta$ denote a primitive $2\mu$-th root of unity, and let $ K =\QQ(\zeta)$. Then for all $r\geq 0$ there exists a representation $\varepsilon_r \colon \mathfrak{t}_\mu \to \GL_{r+1}(K) $    
    if
        \begin{eqnarray*}
            \varepsilon_r(u) = 
            \begin{cases}
                \exp((\zeta+\zeta^{-1})A_r) \quad & \text{if } u = T, \\
                d_r^{\text{y}}((-1)^j) & \text{if }u = S.
            \end{cases}
        \end{eqnarray*}
    holds. In fact, for $ r\geq 1 $ we have
        \begin{eqnarray*}
            \varepsilon_r \cong \Sym^{r}\varepsilon_1
        \end{eqnarray*}
    is true. 
\end{thm}

\begin{proof}
    Utilizing $ \varepsilon_r$ we will induct over $r$.  
    
    When $r=0$ this is the trivial representation. 
    
    Now set $r=1$, so that $\varepsilon_1(S)^2 = -I$ and from which we find
    \[
    \varepsilon_1(T) = \exp\!\left((\zeta+\zeta^{-1})\left(\begin{smallmatrix}
        0&0\\
        1&0
    \end{smallmatrix}\right)\right) = \left(\begin{matrix}
        1&0\\\
        \zeta+\zeta^{-1}&1
    \end{matrix}\right),
    \]
    holds. Thus,
    \[
    \varepsilon_1(R) = \varepsilon_1(S^{-1})\varepsilon_1(T) = \left(\begin{matrix}
        0&1\\
        -1&0
    \end{matrix}\right)\left(\begin{matrix}
        1&0\\\
        \zeta+\zeta^{-1}&1
    \end{matrix}\right) = \left(\begin{matrix}
        \zeta+\zeta^{-1}&1\\
        -1&0
    \end{matrix}\right)
    \]
    is satisfied. We see that the characteristic polynomial of $\varepsilon_1(R)$ is 
        \begin{eqnarray*}
            X^2-(\zeta+\zeta^{-1})X+1
        \end{eqnarray*}
    so that $\varepsilon_1(R)$ is diagonalizable with eigenvalues $\zeta$ and $\zeta^{-1}$. Since $ \zeta^\mu =\zeta^{-\mu}=-1$, we find that $\varepsilon_1(R)^\mu$ is diagonalizable with both eigenvalues equal to $-1$, and hence $\varepsilon_1(R)^\mu = -I = \varepsilon_1(S)^2$. This proves the claim when $r=1$.

    Now for $r\geq 2$ it suffices to show that there exists a basis for $\Sym^r\varepsilon_1$ such that $ S$ and $ T $ act via the matrices $\varepsilon_r(S)$ and $\varepsilon_r(T)$; the identity $\varepsilon_r(R)^\mu = \varepsilon_r(S)^2$ then follows immediately from $\Sym^r\varepsilon_1 $ being a group homomorphism and the $r=1$ case. 
    
    Recall that if we let the standard basis for $\Sym^1\varepsilon_1 = \varepsilon_1$ be $e$ and $f$ where
        \begin{eqnarray*}
            e := \left(
                \begin{matrix}
                    1\\0
                \end{matrix}\right) 
            \quad \text{and} \quad
            f := \left(
                \begin{matrix}
                    0\\1
                \end{matrix}\right),
    \end{eqnarray*}
    then the vector space underlying $\Sym^r\varepsilon_1$ is spanned by monomials $e^if^j$ with $i+j=r$, and $\gamma\in \mathfrak{t}_\mu $ acts by
        \begin{equation}\label{eq:sympower}
            \left(\Sym^r\varepsilon_1(\gamma)\right) \cdot e^if^j = (\varepsilon_1(\gamma)e)^i(\varepsilon_1(\gamma)f)^j.    
        \end{equation}
    Now consider the rescaled basis $ w_k = \frac{e^k f^{r-k}}{k!\,(r-k)!}$ for $ k=0,1,\ldots,r $.

    With respect to our alternative basis, we now describe the {action of $ S $.} Since $ \varepsilon_1(S)\cdot e = f $ and $ \varepsilon_1(S) \cdot f = -e $, equation \eqref{eq:sympower} gives
        \begin{eqnarray*}
            \left(\Sym^r\varepsilon_1(S)\right)\cdot w_k = \frac{f^k\cdot(-e)^{r-k}}{k!\,(r-k)!} = (-1)^{r-k}\,w_{r-k}.
        \end{eqnarray*}
    Hence in the basis $(w_0, w_1, \ldots, w_r)$ the matrix of $\Sym^r\varepsilon_1(S)$ is anti-diagonal with $(k,r-k)$-entry equal to $(-1)^{r-k}$. The upper-right entry (at position $(0,r)$) equals $(-1)^r$, and the signs alternate along the anti-diagonal, exactly as described for $\varepsilon_r(S)$.

    As for the action of $ T $, since $\varepsilon_1(T)\cdot e = e + (\zeta+\zeta^{-1})f$ and $\varepsilon_1(T)\cdot f = f$, equation \eqref{eq:sympower} gives
    \begin{align*}
    \left(\Sym^r\varepsilon_1(T)\right)\cdot w_k
    &= \frac{\bigl(e+(\zeta+\zeta^{-1})f\bigr)^k f^{r-k}}{k!\,(r-k)!} \\
    &= \frac{1}{k!\,(r-k)!}\sum_{m=0}^{k}\binom{k}{m}(\zeta+\zeta^{-1})^m e^{k-m}f^{r-k+m}.
    \end{align*}
    Setting $\ell = k-m$ (so the monomial is $e^\ell f^{r-\ell}$), this becomes
    \[
    \sum_{\ell=0}^{k} \frac{(\zeta+\zeta^{-1})^{k-\ell}}{(k-\ell)!\,(r-k)!}\cdot \frac{e^\ell f^{r-\ell}}{\ell!\,(r-\ell)!}\cdot \ell!\,(r-\ell)!
    = \sum_{\ell=0}^{k} \frac{\ell!\,(\zeta+\zeta^{-1})^{k-\ell}}{\ell!\,(k-\ell)!} \cdot \frac{(r-\ell)!}{(r-k)!} \cdot w_\ell,
    \]
    which simplifies to
    \[
    \left(\Sym^r\varepsilon_1(T)\right)\cdot w_k = \sum_{\ell=0}^{k} \frac{\ell!\,(\zeta+\zeta^{-1})^{k-\ell}}{(k-\ell)!}\, w_\ell.
    \]
    Thus, in the basis $(w_k)$, the $(\ell,k)$-entry of $\Sym^r\varepsilon_1(T)$ (for $\ell \leq k$) is
    \[
    \frac{\ell!\,(\zeta+\zeta^{-1})^{k-\ell}}{(k-\ell)!}.
    \]
    On the other hand, since $(A_r)_{i,j}= i\,\delta_{i,j+1}$, we compute $(A_r^m)_{i,j} = \frac{i!}{(i-m)!}\,\delta_{i,j+m}$, and hence
    \[
    \bigl[\exp\!\left((\zeta+\zeta^{-1})A_r\right)\bigr]_{\ell,k} = \frac{(\zeta+\zeta^{-1})^{\ell-k}}{(\ell-k)!}\cdot\frac{\ell!}{k!}\Bigg|_{\substack{\text{lower}\\\text{triangular}}}
    \]
    Thus, $A_r$ is sub-diagonal and so $\exp(A_r)$ is lower-triangular with $(\ell,k)$-entry (for $\ell \geq k$) equal to
    \[
    \frac{(\zeta+\zeta^{-1})^{\ell-k}}{(\ell-k)!}\cdot\frac{\ell!}{k!}.
    \]
    Renaming $(\ell, k) \mapsto (k, \ell)$ to match the $\Sym^r$ computation above, we see the two expressions agree:
    \[
    \frac{\ell!\,(\zeta+\zeta^{-1})^{k-\ell}}{(k-\ell)!\,} = \frac{(\zeta+\zeta^{-1})^{k-\ell}}{(k-\ell)!}\cdot \ell!,
    \]
    This shows that in the $w$-basis the $(\ell, k)$-entry of $\Sym^r\varepsilon_1(T)$ equals the $(\ell,k)$-entry of $\exp((\zeta+\zeta^{-1})A_r)$, completing the verification that $\varepsilon_r(T) = \Sym^r\varepsilon_1(T)$ in the $w$-basis. This completes the proof.
\end{proof}

\end{document}
